Tables~\ref{tab:generation}--\ref{tab:raw-generation} distinguish raw P decoding, separately trained instruction responses, and label-in-text references. All values below are class-macro scores on held-out templates at 6000 G steps.

For IMU, label success is 0.838 [0.737, 0.936] (correct-minus-different-class 0.832 [0.731, 0.933]); sentence success is 0.689 [0.198, 0.971] (correct-minus-different-class 0.686 [0.198, 0.966]). Per-seed label scores are 0.928, 0.740, 0.847; sentence scores are 0.916, 0.193, 0.959, in seed order 42, 1042, 2042.

For Audio, label success is 0.056 [0.000, 0.161] (correct-minus-different-class 0.051 [0.000, 0.144]); sentence success is 0.064 [0.000, 0.186] (correct-minus-different-class 0.060 [0.000, 0.167]). Per-seed label scores are 0.000, 0.000, 0.169; sentence scores are 0.000, 0.000, 0.192, in seed order 42, 1042, 2042. Two of the three seeds have zero exact success in both families.

For held-out rule prompts, the differences in macro-averaged joint success (correct sensor minus a sensor from another class) are IMU $-2.1\times10^{-4}$ [-0.007, 0.004]; Audio -0.001 [-0.007, 0.006]. Requiring both reversed mappings to be correct for the same sensor and template gives the secondary scores IMU 0.005 [0.001, 0.013]; Audio 0.001 [0.000, 0.002]. Reversing the output mapping while keeping the sensor fixed tests whether responses follow the supplied rule as well as the sensor. These deterministic templates cannot demonstrate general instruction following, and sensor-dependent labels or fixed sentences alone do not show reliable rule application. The label-in-text references are practical comparisons with a different input interface; their architectures are not matched to the bridge. The evidence includes every response, including invalid and repetitive output.
